\documentclass[12pt]{article}

\usepackage[spanish]{babel}
\usepackage[utf8]{inputenc}
\usepackage[T1]{fontenc}
\usepackage{amsmath, amssymb}
\usepackage{enumitem}
\usepackage[a4paper, margin=2.5cm]{geometry}
\usepackage{lmodern}

\begin{document}
	
	\begin{center}
		{\Large \textbf{Guía de Ejercicios -- Evaluación II}}\\[0.3cm]
		{\large \textbf{Contextualizada al área de Agronomía}}
	\end{center}
	
	\vspace{0.4cm}
	
	\noindent
	\textbf{Instrucciones:} En cada caso realice todo el desarrollo ordenado y responda claramente la pregunta planteada.
	
	\vspace{0.5cm}
	
	\begin{enumerate}
		
		% =====================================================
		\item \textbf{Producción agrícola en distintos sectores}
		
		De acuerdo con el registro de producción de una cooperativa agrícola, se tiene la siguiente información sobre la cantidad de sacos cosechados en distintos sectores:
		
		\[
		\begin{array}{|l|c|}
			\hline
			\textbf{Sector de cultivo} & \textbf{Número de sacos} \\
			\hline
			\text{Maíz} & 125 \\
			\text{Trigo} & 98 \\
			\text{Papas} & 143 \\
			\text{Tomates} & 76 \\
			\hline
		\end{array}
		\]
		
		\begin{enumerate}[label=\alph*)]
			\item Determine el número total de sacos cosechados.
			\item Determine la diferencia entre la producción de maíz y la producción de tomates.
		\end{enumerate}
		
		\vspace{0.4cm}
		
		% =====================================================
		\item \textbf{Aplicación de fertilizante}
		
		Para fertilizar una plantación, se deben tratar 
		\[
		(3n^2+5n)
		\]
		plantas, y la cantidad de fertilizante líquido que se aplica a cada una es de
		\[
		(2n+7)\ \text{mL}.
		\]
		
		Determine la cantidad total de fertilizante, en mL, que se debe utilizar.
		
		\vspace{0.4cm}
		
		% =====================================================
		\item \textbf{Duración de un tratamiento fitosanitario}
		
		En un cultivo afectado por una plaga, un agrónomo dispone de \(x\) envases de un producto fitosanitario. Cada envase contiene
		\[
		(x^2-x+72)
		\]
		gramos del producto.
		
		Si diariamente se aplican
		\[
		(x^2-9x)
		\]
		gramos, ¿cuántos días dura el tratamiento?
		
		\vspace{0.4cm}
		
		% =====================================================
		\item \textbf{Costo total de insumos agrícolas}
		
		Un productor compra una cantidad de insumos agrícolas dada por:
		\[
		Q=\frac{x^2-5x+6}{x-3}
		\qquad \text{[unidades]}
		\]
		
		y cada unidad tiene un valor de:
		\[
		P=\frac{x^3+x^2+4x+4}{x^2-x-2}
		\qquad \text{[pesos por unidad]}
		\]
		
		El costo total \(C\) de la compra está dado por:
		\[
		C=P\cdot Q
		\]
		
		\begin{enumerate}[label=\alph*)]
			\item Simplifique el costo total a su mínima expresión.
			\item Evalúe el costo total para \(x=1500\).
		\end{enumerate}
		
		\vspace{0.4cm}
		
		% =====================================================
		\item \textbf{Modelo algebraico de concentración de nutrientes}
		
		Un laboratorio agronómico estudia la concentración de nutrientes presentes en una solución de riego. La expresión obtenida es:
		
		\[
		C=
		\left[
		\frac{x^2-x-12}{x^2-9}
		\cdot
		\frac{x^3+x^2-2x}{x^2-2x-8}
		\div
		\frac{x-1}{2x^2-6x}
		\right]
		\]
		
		Aplicando factorización y simplificación, encuentre la expresión más reducida que representa la concentración \(C\).
		
		\vspace{0.4cm}
		
		% =====================================================
		\item \textbf{Simplificación de expresiones algebraicas en contexto agrícola}
		
		Simplifique las siguientes expresiones:
		
		\begin{enumerate}[label=\alph*)]
			\item 
			\[
			\frac{[(8x^2-5)+4]^2}{(2x+7)^2-(4x-7)}
			\]
			
			\item
			\[
			\frac{(3x-2)(3x+4)}{(5x+3)(3+5x)-9}
			\]
		\end{enumerate}
		
	\end{enumerate}
	
\end{document}